\uuid{y1Xv}
\titre{Espace des polyômes}
\auteur{hanani ; mzari}
\organisation{exo7}
\licence{CC BY-NC-SA 4.0}
\datecreate{2020-01-20}
\niveau{L1}
\module{Algèbre}
\chapitre{Espace vectoriel}
\sousChapitre{Définition, sous-espace}
\difficulte{3}
\theme{Espaces vectoriels, Espace des polyômes}
\exo7qcmid{lille-2.espaces-vectoriels.013}

\contenu{
\begin{question}
\qclassification{106.01}
\qtags{motcle=Espace des polyômes}
Soit $\Rr_2[X]$ l'espace des polynômes à coefficients réels, de degré inférieur ou égal à $2$, muni des opérations usuelles et
$E=\{P \in \Rr_2[X] \, ; \; P(1)=1\}$. Quelles sont les assertions vraies ?
\begin{answers}
\bad{$E$ est vide.}
\bad{$E$ est stable par addition.}
\good{$E$ n'est pas stable par multiplication par un scalaire.}
\good{$E$ n'est pas un espace vectoriel.}
\end{answers}
\begin{explanations} $E$ n'est pas un sous-espace vectoriel de $\Rr_2[X]$, puisque le polynôme nul n'appartient pas à $ E$.
$E$ n'est stable ni par addition ni par multiplication par un scalaire.
\end{explanations}
\end{question}
}
